\answer{ex:05:1}{$\ell\approx17\um$; $\approx100$일. 뇌 전체로의 배포는 전선의 분지가 맡고, 확산은 각 방출 지점을 $\sim20\um$ 정도로 번지게 할 뿐이다.}
\answer{ex:05:2}{$r=ab/(1+G)$, $S=1/(1+G)$는 정확하다. $+20\,\%$가 아니라 $+1.7\,\%$.}
\answer{ex:05:3}{$G=2$에서 $T^*=1.21\,\text{s}$, $G=9$에서 $0.19\,\text{s}$. 현실적으로는 $G\sim2$--$4$, 억제는 $3$--$5$배.}
\answer{ex:05:4}{전체 발화율 $\lambda$에 대해 $\mathrm{CV}=1/\sqrt{2\lambda\tau_c}\approx9\cdot10^{-4}$. 뉴런 하나라면 $\tau_c=12.5\,\text{s}$가 필요하다.}
\answer{ex:05:5}{두 경우 모두 $c\approx0.40$; $\mathrm{CV}=0.050$ 대 $0.158$, $\sqrt{10}$배. 개별 뉴런의 발화율은 여전히 $a_i$에 비례한다. 조절되는 것은 합뿐이다.}
\answer{ex:05:6}{$N_1=\overline{A\vee B}$, $N_2=\overline{A\vee N_1}$, $N_3=\overline{B\vee N_1}$, $N_4=\overline{N_2\vee N_3}$, XOR $=N_5=\overline{N_4}$. 전환 한 번에 $\tau_c\ln2$가 걸리므로 게이트 $4$개를 거치는 경로는 XOR 하나에 $2.8\,\text{s}$.}
\answer{ex:05:7}{(가)~$\tau_\gamma=6/\ln3\approx5.5\,\text{s}$. (나)~$\bar D<4\,\text{s}$ 대 $D<2.2\,\text{s}$: 지연을 예측할 수 없으면 더 참을성이 생긴다.}
\answer{ex:05:8}{$k_2=1/\tau_d=10\,\text{s}^{-1}$, $k_1=1/\tau_d-1/\tau_\gamma=9.5\,\text{s}^{-1}$; $\tau_\gamma=1/(k_2-mk_1)$. $+2.6\,\%$에서 두 배가 된다($+5.3\,\%$에서는 지평이 무한대).}
